% Generated by build_manuscript.py from 01_introduction.md
\hypertarget{introduction}{%
\section{1 Introduction}\label{introduction}}

Estimating the remaining compressive strength of an impacted composite requires connecting observable damage to a mechanical assessment. Surface appearances, internal damage images and compression-after-impact (CAI) strength measurements provide complementary information about that problem. Public composite-impact datasets associate these observations with corresponding specimens \citep{hasebe2022data, hasebe2025data}. They support image-based strength assessment and a further question: which internal observations should be acquired when only part of the image will be available? The value of an observation then depends on its contribution to an evolving strength estimate.

C-scan image regression establishes a route from internal damage information to mechanical assessment. Recent work directly predicts CAI strength and impact energy from C-scan images using a convolutional model \citep{mack2026}. Selective acquisition adds a decision before that mapping can be applied to complete input. A policy must choose which missing region to observe next using evidence already available. Two sequences can eventually acquire the same subset yet provide different intermediate predictions because useful information becomes available at different stages. Assessing only the final estimate would overlook this distinction between observation content and availability.

Adaptive ultrasonic inspection provides a precedent for choosing subsequent measurements from earlier observations \citep{fuentes2020}. Dynamic feature-acquisition methods more generally connect information requests with prediction utility and acquisition cost \citep{shim2018, janisch2019, covert2023}. For CAI assessment, the relevant utility is how the acquired evidence changes a strength estimate. Visually conspicuous damage may guide an initial observation, but its predictive value must be assessed together with other available information. This motivates a sequential objective that accounts for prediction error throughout acquisition as well as at the endpoint.

The modalities supply different information at different stages. A surface image is available initially, whereas internal content enters the state only after its region has been acquired. Surface cues can therefore initialize a search, and subsequent internal observations can change both the predicted strength and the next acquisition decision. Implementing this progression requires an explicit state interface. The policy must distinguish observed from unobserved regions and relate each permitted descriptor to its location, acquisition history and remaining budget. This interface connects evidence availability to the task objective rather than treating the next region as a fixed function of surface appearance alone.

We develop a state-dependent multimodal acquisition method with three distinct roles. A frozen vision-language model (VLM) supplies surface-region priors, a learned actor selects each internal cell, and a separate frozen predictor estimates CAI strength from the current observation set. The actor is trained using normalized trajectory error area together with terminal absolute error. At each acquisition step, the new internal descriptor and updated prediction enter its state, while the model parameters remain fixed. A common predictor across strategies holds the assessment mapping constant, enabling comparison of the subsets supplied and the timing of their availability.

The study makes three contributions. First, it formulates sequential C-scan acquisition around CAI prediction quality over the acquisition process, rather than only complete-image regression. Second, it connects surface priors, observed internal descriptors and the current strength estimate to subsequent acquisition through a state-dependent policy. Third, it characterizes this policy under a common predictor through same-cap quality, signed timing contributions and matched-quality acquisition requirements. Offline replay on 50 validation specimens compares shared, surface-conditioned and feedback-driven orders. The analysis asks where the resulting prediction value is accumulated and how that process relates to the cost of attaining an empirical quality target; the selection and evaluation scope is specified in Section 4.
